\section{Experimental Setup}

We design experiments to test the foundational assumption of Progressive Gradient Orthogonalization: that early training gradients primarily capture spurious feature directions.

\subsection{Research Questions}

\textbf{RQ1:} Does the accumulated gradient subspace align more strongly with spurious (background) directions than core (bird type) directions during early training?

\textbf{RQ2:} Is the alignment difference measurable given realistic subspace ranks in high-dimensional parameter spaces?

\subsection{Dataset}

\textbf{Waterbirds} \citep{sagawa2020distributionally}: A spurious correlation benchmark where background type (water/land) correlates 95\% with bird type (waterbird/landbird) in training data.

\begin{table}[h]
\centering
\caption{Waterbirds dataset statistics.}
\begin{tabular}{@{}lccc@{}}
\toprule
Split & Samples & Groups & Spurious Correlation \\
\midrule
Train & 4,795 & 4 & 95\% \\
Validation & 1,199 & 4 & 95\% \\
Test & 5,794 & 4 & Balanced \\
\bottomrule
\end{tabular}
\end{table}

\subsection{Model and Training}

\textbf{Model:} ResNet-50 pretrained on ImageNet, final fully-connected layer replaced with a 2-class linear head (25.6M trainable parameters).

\begin{table}[h]
\centering
\caption{Training hyperparameters.}
\begin{tabular}{@{}ll@{}}
\toprule
Parameter & Value \\
\midrule
Optimizer & SGD (momentum=0.9, weight\_decay=1e-4) \\
Learning rate & 1e-3 (step decay at epochs 60, 75) \\
Batch size & 128 \\
Total epochs & 90 \\
\bottomrule
\end{tabular}
\end{table}

\textbf{Gradient Accumulation:}
\begin{itemize}
    \item Accumulation window: epochs 1-10
    \item Samples accumulated: 1 gradient per epoch (last batch)
    \item SVD rank: $k=50$
\end{itemize}

\subsection{Evaluation Metrics}

\begin{itemize}
    \item Spurious alignment $> 0.70$ at epoch 10
    \item Core alignment $< 0.30$ at epoch 10
    \item Measurement epochs: 5, 10, 45
\end{itemize}
